\begin{abstract}
\noindent The USPTO publishes every US trademark application filed since
1884, together with the history of what happened to it. We re-parse all
13.99 million of these case files into 242 million dated events, so that
every mark can be followed step by step from filing through examination,
registration, the sworn five-year proof of continued use, and the
year-ten renewal, with counsel of record, filing basis and owner domicile
read from the record, and owners linked by name to SEC reporting,
Form~D fundraising rounds and patent assignees. The dates matter because
the USPTO's terminal status codes run the maintenance deadlines together:
a registration cancelled for missing the year-six proof of use and one
cancelled at the year-ten renewal share a code. Once each cancellation is
dated, the record shows which deadline a lapsed mark failed, and the
five-year proof becomes something the literature has not had: a dated,
product-level record of whether the goods behind a registered mark were
still being sold. On this corpus we build a measure of where a filing's
language sits relative to its industry. A topic model reduces each
goods/services description to a mix of themes---an online offering, say,
or an artificial-intelligence application---and we compare that mix, by
Kullback--Leibler divergence, with what the filing's own Nice class filed
in the five years before its date and in the five years after. The
average of the two comparisons is the filing's \emph{atypicality}: how
unusual its language is for its industry. Their difference is its
\emph{lead}: whether its language ran ahead of or behind where the
industry's language was moving. We validate the measure by rescoring
identical filings under three partitions of the vocabulary, two reference
weightings, five window constructions, a second model seed, and with the
pervasive verbatim text duplication removed. Lead survives every
variation and predicts cancellation for non-use at the five-year proof
($+2.3$, $+2.6$ and $+1.9$ points under the three partitions; $+1.4$ with
design controls, $t = 8.3$). Atypicality does not: its apparent effects
belong to the partition or to a control. We document two hazards for
trademark-text research---scoring words rather than themes measures
description length ($r = -0.68$ with log distinct terms), and firm-level
correlations built on word scoring dissolve or reverse under theme
scoring on identical samples. Within firm, the measure is orthogonal to
patenting. The corpus, the scores, the code, and an interactive per-class
appendix are public.

\medskip
\noindent\emph{Keywords:} trademarks; administrative data; text as data;
topic models; Kullback--Leibler divergence; innovation indicators.\\
\emph{JEL:} C81, O34, O31, L10.
\end{abstract}
